\section{The \protect\Lovasz{} Local Lemma}

% [CORRECTED] was: "$\Pr\of{B_i} \leq 1$" --- that holds of every event whatsoever, so it is no hypothesis, and
% $B_1 = \Omega$ refutes the sentence after it. With $\Pr\of{B_i} < 1$ the claim is not just true but exact:
% $\Pr\of{\bigcap B_i^C} = \prod\parenth{1 - \Pr\of{B_i}}$, which is positive precisely when every factor is.
Suppose $B_1, \ldots, B_n$ are independent events with $\Pr\of{B_i} < 1$. It's pretty easy to see that $\Pr\of{\bigcap B_i^C} > 0$. Essentially, we think of the $B_i$ as \textit{bad events}: events we \textit{don't} want to see happen.

We want to extend this. In what way exactly? Here's a motivating question.

In a $k$-uniform hypergraph where each edge intersects at most $D$ others, what condition on $k$ makes the hypergraph $2$-colourable?

What we'd want to do is randomly colour vertices: each edge $f$ is monochromatic with probability $\Pr\of{B_f} = \frac{1}{2^{k-1}}$. We would then like to apply some version of the fact that $\Pr\of{\bigcap B_i^C} > 0$.

Two edges sharing a vertex are plainly not independent, so what we ask instead is that each bad event depend on only a few of the others---and a graph is the natural place to record which.

% A [SUSPECT] stood here from 2026-09-23 to 2026-10-07: "$A$ is independent of $B_1, \ldots, B_n$" was not a notion
% these notes defined, and read pairwise it leaves the local lemma false (three events $\set{\xi_i = 1}$ with
% $\xi_3 = 1 \oplus \xi_1 \oplus \xi_2$ are pairwise independent, so the empty graph would be a dependency graph and
% $x_{A_i} = \frac{1}{2}$ would satisfy the hypothesis, yet the three complements cannot all hold). Settled by
% adding the third form of independence to \Cref{Ch1:Def:Independence}, which is the reading the proof below uses.
\begin{boxdefinition}[Dependency Graph] \label{Ch5:Def:Dependency_Graph}
    Given a collection $\mathcal{A}$ of events, a \textbf{dependency graph on $\mathcal{A}$} is a graph with vertices $\setst{v_A}{A \in \mathcal{A}}$ such that $v_A \not\sim v_{B_1}, \ldots, v_{B_n}$ implies that $A$ is independent of the family $B_1, \ldots, B_n$ (in the sense of \Cref{Ch1:Def:Independence}).
\end{boxdefinition}

% [CORRECTED] was: "Another way of phrasing this is that adjacency represents dependence (ie, lack of
% independence)." --- that is the converse of the definition above, which asks that dependence force adjacency and
% says nothing about adjacent events having to depend on each other.
Another way of phrasing this is that dependence forces adjacency: a dependency graph may join more than it has to, but it may not leave a dependence unrecorded.

% [CORRECTED] was: "Borrowing from our use of $\sim$ for adjacency, we will denote that events $A$ and $B$ are
% dependent (ie, not independent) by $A \sim B$." --- read that way, $\sim$ is pairwise dependence, and the local
% lemma below is then false. For $\xi_1, \xi_2$ independent fair bits and $\xi_3 = 1 \oplus \xi_1 \oplus \xi_2$, the
% events $A_i = \set{\xi_i = 1}$ are pairwise independent, so every product below is empty; $x_{A_i} = \frac{1}{2}$
% satisfies the hypothesis, yet $A_1^C \cap A_2^C \cap A_3^C = \emptyset$. What the lemma needs is adjacency in a
% dependency graph, where a non-neighbour is independent of $A$ as a family and not one at a time.
\begin{boxlnotation}[Dependence]
    Fix a dependency graph on $\mathcal{A}$. Borrowing from our use of $\sim$ for adjacency, we will write $A \sim B$ for $v_A \sim v_B$, and read it as saying that $A$ and $B$ are permitted to depend on each other.
\end{boxlnotation}

The extension we were after weighs each event against the ones it depends on.

\begin{boxtheorem}[\Lovasz{} Local Lemma] \label{Ch5:Thm:LLL}
    % [CORRECTED] was: "$0 < x_A \leq 1$" --- at $x_A = 1$ on an event with no neighbours the hypothesis reads
    % $\Pr(A) \leq 1$, which is no hypothesis at all, and $A = \Omega$ then refutes the conclusion. Strictness is
    % load-bearing twice in the proof: it keeps $\prod_{B \sim A}\parenth{1 - x_B}$ nonzero, so the conditional
    % probabilities are defined, and it is the whole of the final $\prod_{A}\parenth{1 - x_A} > 0$.
    % [CORRECTED] was: "a collection $\mathcal{A}$ of events" --- the conclusion needs $\mathcal{A}$ finite. On
    % $\set{0, 1}^{\N}$ with fair coins, $A_n = \setst{\omega}{\omega_n = 1}$ are independent, so every product is
    % empty and $x_{A_n} = \frac{1}{2}$ satisfies the hypothesis; but $\bigcap A_n^C$ is one sequence, of
    % probability $0$. Finiteness is spent in the same last step as strictness, on $\prod_A \parenth{1 - x_A} > 0$.
    Given a finite collection $\mathcal{A}$ of events and real numbers $0 < x_A < 1$ such that for all $A \in \mathcal{A}$,
    \begin{align*}
        \Pr\of{A} \leq x_A \prod_{B \sim A} \parenth{1 - x_B}
    \end{align*}
    we can conclude that
    \begin{align*}
        \Pr\of{\bigcap_{A \in \mathcal{A}} A^C} > 0
    \end{align*}
\end{boxtheorem}
\begin{proof}
    % [FILLED] proof supplied on 2026-10-07: the lecture of 16 September left it for ``next time'', and no later lecture in these notes gives it
    We show that for every $A \in \mathcal{A}$ and every $S \ssq \mathcal{A} \setminus \set{A}$,
    \begin{align*}
        \Pr\of{A \;\middle\vert\; \bigcap_{B \in S} B^C} \leq x_A
    \end{align*}
    by induction on $\abs{S}$. (Every conditional probability below is defined: the last display of the proof, applied to $S$ in place of $\mathcal{A}$, gives $\Pr\of{\bigcap_{B \in S} B^C} \geq \prod_{B \in S} \parenth{1 - x_B} > 0$, and only uses the claim for smaller sets.) When $S = \emptyset$ this is the hypothesis, as $\prod_{B \sim A} \parenth{1 - x_B} \leq 1$. Otherwise split $S$ into $S_1 = \setst{B \in S}{B \sim A}$ and $S_2 = S \setminus S_1$, and write $E_i = \bigcap_{B \in S_i} B^C$, so that
    \begin{align*}
        \Pr\of{A \mid E_1 \cap E_2} = \frac{\Pr\of{A \cap E_1 \mid E_2}}{\Pr\of{E_1 \mid E_2}}
    \end{align*}
    The numerator is at most $\Pr\of{A \mid E_2}$, and $E_2$ is built out of events not adjacent to $A$, so this is $\Pr\of{A} \leq x_A \prod_{B \sim A} \parenth{1 - x_B}$. For the denominator, enumerate $S_1 = \set{B_1, \ldots, B_r}$ and write
    \begin{align*}
        \Pr\of{E_1 \mid E_2} = \prod_{i = 1}^{r} \Pr\of{B_i^C \;\middle\vert\; B_1^C \cap \cdots \cap B_{i - 1}^C \cap E_2} \geq \prod_{i = 1}^{r} \parenth{1 - x_{B_i}} \geq \prod_{B \sim A} \parenth{1 - x_B}
    \end{align*}
    where each factor is bounded by the induction hypothesis, since each conditions on fewer than $\abs{S}$ events. Dividing, $\Pr\of{A \mid E_1 \cap E_2} \leq x_A$.

    Finally, enumerate $\mathcal{A} = \set{A_1, \ldots, A_n}$, so that
    \begin{align*}
        \Pr\of{\bigcap_{i = 1}^{n} A_i^C} = \prod_{i = 1}^{n} \Pr\of{A_i^C \;\middle\vert\; A_1^C \cap \cdots \cap A_{i - 1}^C} \geq \prod_{i = 1}^{n} \parenth{1 - x_{A_i}} > 0
    \end{align*}
\end{proof}
